\section{スキャンマッチング残差に関するヤコビアンの導出}
\label{sec:scan_matching}

本節では，式(\ref{eq:scan_matching_residual})に示すスキャンマッチングに関するヤコビアンの導出を行う．
なお表記の簡略化のために，${}^{L}{\bf r}(t) = {}^{M}{\bf p} - \left(R(t) {}^{I}{\bf p} + {\bf p}(t) \right)$とし，${}^{L}r(t) = {\bf n}^{\top} {}^{L}{\bf r}(t)$と表記する．
今回の目的は，式(\ref{eq:scan_matching_residual})に示す残差の各制御点$\left( R_{j}, {\bf p}_{j} \right)$に対するヤコビアン$\partial {}^{L}r(t) / \partial R_{j}$，$\partial {}^{L}r(t) / \partial {\bf p}_{j}$を求めることである．
式(\ref{eq:continuous_rotation_matrix})，(\ref{eq:continuous_translation_vector})に示す通り，$\left( R(t), {\bf p}(t) \right)$は制御点$\left( R_{j}, {\bf p}_{j} \right)$を用いて表される．
そのため連鎖則を用いて，$\partial {}^{L}r(t) / \partial R_{j}$，$\partial {}^{L}r(t) / \partial {\bf p}_{j}$は以下のように計算される．
%
\begin{align}
  \frac{ \partial {}^{L}r(t) }{ \partial R_{j} }
  =
  \frac{ \partial {}^{L}r(t) }{ \partial R(t) }
  \frac{ \partial R(t) }{ \partial R_{j} }  
\end{align}
%
\begin{align}
  \frac{ \partial {}^{L}r(t) }{ \partial {\bf p}_{j} }
  =
  \frac{ \partial {}^{L}r(t) }{ \partial {\bf p}(t) }
  \frac{ \partial {\bf p}(t) }{ \partial {\bf p}_{j} }
\end{align}
%
以下，まず並進ベクトル${\bf p}_{j}$に関するヤコビアンについて考える．




\subsection{並進に関するヤコビアンの導出}

まず$\partial {}^{L}r(t) / \partial {\bf p}(t)$は以下の様に変形できる．
%
\begin{align}
  \frac{ \partial {}^{L}r(t) }{ \partial {\bf p}(t) }
  =
  \frac{ \partial {}^{L}r(t) }{ \partial {}^{L}{\bf r}(t) }
  \frac{ \partial {}^{L}{\bf r}(t) }{ \partial {\bf p}(t) }
  = {\bf n}^{\top} \left( -I \right)
  = -{\bf n}^{\top}
\end{align}
%
これは通常のスキャンマッチングを考えた際のヤコビアンと同じため導出の詳細は省く．

次に$\partial {\bf p}(t) / \partial {\bf p}_{j}$ついて考える．
ここでまず，式(\ref{eq:continuous_translation_vector})を以下の様に書き直す．
%
\begin{align}
  \begin{gathered}
  {\bf p}(t) = {\bf p}_{0} + \sum_{j=1}^{3} \lambda_{j} {\bf d}_{j} \\
  %
  {\bf d}_{j} = {\bf p}_{j} - {\bf p}_{j-1}
  \end{gathered}
  \label{eq:continuous_translation_vector_delta}
\end{align}
%
すなわち${\bf p}(t)$は，以下の様に${\bf p}_{0}$と${\bf d}_{j}$に関する関数として記述することができる．
%
\begin{align}
  {\bf p} \left(
    {\bf p}_{0},
    {\bf d}_{1} \left( {\bf p}_{0}, {\bf p}_{1} \right),
    {\bf d}_{2} \left( {\bf p}_{1}, {\bf p}_{2} \right),
    {\bf d}_{3} \left( {\bf p}_{2}, {\bf p}_{3} \right)
  \right)
\end{align}
%
$j=1,2$の場合，${\bf p}_{j}$は${\bf d}_{j}$と${\bf d}_{j+1}$に依存する．
そのため$\partial {\bf p}(t) / \partial {\bf p}_{j}$は，連鎖則を用いて以下の様に計算できる．
%
\begin{align}
  \frac{ \partial {\bf p}(t) }{ \partial {\bf p}_{j} }
  =
  \frac{ \partial {\bf p}(t) }{ \partial {\bf d}_{j} }
  \frac{ \partial {\bf d}_{j} }{ \partial {\bf p}_{j} }
  +
  \frac{ \partial {\bf p}(t) }{ \partial {\bf d}_{j+1} }
  \frac{ \partial {\bf d}_{j+1} }{ \partial {\bf p}_{j} }
  ~~~(j=1,2)
\end{align}
%
$j=0$と$j=3$の場合は以下となる．
%
\begin{align}
  \frac{ \partial {\bf p}(t) }{ \partial {\bf p}_{0} }
  =
  \frac{ \partial {\bf p}(t) }{ \partial {\bf p}_{0} }
  +
  \frac{ \partial {\bf p}(t) }{ \partial {\bf d}_{1} }
  \frac{ \partial {\bf d}_{1} }{ \partial {\bf p}_{0} }
\end{align}
%
\begin{align}
  \frac{ \partial {\bf p}(t) }{ \partial {\bf p}_{3} }
  =
  \frac{ \partial {\bf p}(t) }{ \partial {\bf d}_{3} }
  \frac{ \partial {\bf d}_{3} }{ \partial {\bf p}_{3} }
\end{align}
%

式(\ref{eq:continuous_translation_vector})より，それぞれの偏微分の結果は以下となる．
%
\begin{align}
  \begin{gathered}
    \frac{ \partial {\bf p}(t) }{ \partial {\bf p}_{0} }
    =
    \left( 1 - \lambda_{1} \right) I \\
    %
    \frac{ \partial {\bf p}(t) }{ \partial {\bf p}_{1} }
    =
    \left( \lambda_{1} - \lambda_{2} \right) I \\
    %
    \frac{ \partial {\bf p}(t) }{ \partial {\bf p}_{2} }
    =
    \left( \lambda_{2} - \lambda_{3} \right) I \\
    %
    \frac{ \partial {\bf p}(t) }{ \partial {\bf p}_{3} }
    =
    \lambda_{3} I \\
  \end{gathered}
\end{align}
%
これらに$\partial {}^{L}r(t) / \partial {\bf p}(t) = -{\bf n}^{\top}$を左から掛けた行列が，$\partial {}^{L}r(t) / \partial {\bf p}_{j}$となる．

なお式(\ref{eq:b-spline_basis})と式(\ref{eq:b-spline_lambda})を利用すると，$\partial {}^{L}r(t) / \partial {\bf p}_{j} = -\beta_{j} {\bf n}^{\top}$と書くこともできる．
ここで$\sum_{j} \beta_{j} = 1$となることを考えると，$\partial {}^{L}r(t) / \partial {\bf p}_{j}$は，時刻$t$の並進${\bf p}(t)$に関するヤコビアンを，各制御点の寄与度$\beta_{j}$で分配したものと見なすこともできる．
このように，並進は制御点の線形結合で表されるため，そのヤコビアンは比較的容易に求められる．
一方，回転はリー群上の非線形な積として表現されるため，並進のような単純な重み付け分配にはならない．







\subsection{回転に関するヤコビアンの導出}

まず式(\ref{eq:scan_matching_residual})の$R_{j}$に関するヤコビアンは，以下の連鎖則で求められる．
%
\begin{align}
  \frac{ \partial {}^{L}r(t) }{ \partial R_{j} }
  =
  \frac{ \partial {}^{L}r(t) }{ \partial {}^{L}{\bf r}(t) }
  \frac{ \partial {}^{L}{\bf r}(t) }{ \partial R(t) }
  \frac{ \partial R(t) }{ \partial R_{j} }
\end{align}
%
$\partial {}^{L}r(t) / \partial {}^{L}{\bf r}(t)$および$\partial {}^{L}{\bf r}(t) / \partial R(t)$に関して通常のスキャンマッチングでも利用されるため詳細な導出は省くが，計算結果は以下となる．
%
\begin{align}
  \frac{ \partial {}^{L}r(t) }{ \partial {}^{L}{\bf r}(t) }
  \frac{ \partial {}^{L}{\bf r}(t) }{ \partial R(t) }
  =
  {\bf n}^{\top} \left( R(t) {}^{I}{\bf p} \right)^{\wedge}
\end{align}
%

次に$\partial {}^{L}r(t) / \partial R_{j}$について考える．
式(\ref{eq:continuous_translation_vector})に示した様に，$R(t)$も以下の様に記述する．
%
\begin{align}
  \begin{gathered}
  R(t) = R_{0} \prod_{j=1}^{3} {\rm Exp} \left( \lambda_{j} \boldsymbol \phi_{j} \right) \\
  %
  \boldsymbol \phi_{j} = {\rm Log} \left( R_{j-1}^{-1} R_{j} \right)
  \end{gathered}
  \label{eq:continuous_rotation_matrix_delta}
\end{align}
%
すなわち$R(t)$も同様に，$R_{0}$と$\boldsymbol \phi_{j}$に関する関数として記述することができる．
%
\begin{align}
  R \left(
    R_{0},
    \boldsymbol \phi_{1} \left( R_{0}, R_{1} \right),
    \boldsymbol \phi_{2} \left( R_{1}, R_{2} \right),
    \boldsymbol \phi_{3} \left( R_{2}, R_{3} \right)
  \right)
\end{align}
%
ここから，$j=1,2$の場合，$R_{j}$は$\boldsymbol \phi_{j}$と$\boldsymbol \phi_{j+1}$に依存することがわかる．
そのため$j=1,2,$の場合の$\partial R(t) / \partial R_{j}$は，連鎖則を用いて以下のように計算できる．
%
\begin{align}
  \frac{ \partial R(t) }{ \partial R_{j} }
  =
  \frac{ \partial R(t) }{ \partial \boldsymbol \phi_{j} }
  \frac{ \partial \boldsymbol \phi_{j} }{ \partial R_{j} }
  +
  \frac{ \partial R(t) }{ \partial \boldsymbol \phi_{j+1} }
  \frac{ \partial \boldsymbol \phi_{j+1} }{ \partial R_{j} }
\end{align}
%
$j=0$と$j=3$の場合は以下となる．
%
\begin{align}
  \frac{ \partial R(t) }{ \partial R_{0} }
  =
  \frac{ \partial R(t) }{ \partial R_{0} }
  +
  \frac{ \partial R(t) }{ \partial \boldsymbol \phi_{1} }
  \frac{ \partial \boldsymbol \phi_{1} }{ \partial R_{0} }
\end{align}
%
\begin{align}
  \frac{ \partial R(t) }{ \partial R_{0} }
  =
  \frac{ \partial R(t) }{ \partial \boldsymbol \phi_{3} }
  \frac{ \partial \boldsymbol \phi_{3} }{ \partial R_{3} }
\end{align}

ここからまず，$\partial R(t) / \partial R_{0}$について考える．
表記の簡略化のために，$R(t)$を以下のように記述する．
%
\begin{align}
  \begin{gathered}
    R(t) = R_{0} A_{1} A_{2} A_{3} \\
    %
    A_{j} = {\rm Exp} \left( \lambda_{j} \boldsymbol \phi_{j} \right)
  \end{gathered}
\end{align}
%
この表記を用いて，$\partial R(t) / \partial R_{0}$は，以下を計算することで求めることができる．
%
\begin{align}
  \begin{split}
    & \left( {\rm Exp} \left( \delta \boldsymbol \theta \right) R_{0} A_{1} A_{2} A_{3} \right) \left( R_{0} A_{1} A_{2} A_{3} \right)^{-1} \\
  %
  = & {\rm Exp} \left( \delta \boldsymbol \theta \right) R_{0} A_{1} A_{2} A_{3} A_{3}^{-1} A_{2}^{-1} A_{1}^{-1} R_{0}^{-1} \\
  %
  = & {\rm Exp} \left( \delta \boldsymbol \theta \right) \\
  %
  = & I + \left( \delta \boldsymbol \theta \right)^{\wedge} + O \left( \left\| \delta \boldsymbol \theta \right\|^{2} \right)
  \end{split}
\end{align}
%
よって以下を得る．
%
\begin{align}
  \frac{ \partial R(t) }{ \partial R_{0} } = I
\end{align}
%

次に$\partial R(t) / \partial \boldsymbol \phi_{j}$について考える．
これは$j$の値によって計算結果が異なるため一般形式で書くことはできないが，例えば$j=1$の場合は以下のように計算できる．
%
\begin{align}
  \begin{split}
    & R_{0} {\rm Exp} \left( \lambda_{1} \left( \boldsymbol \phi_{1} + \delta \boldsymbol \phi_{1} \right) \right) A_{2} A_{3} \left( R_{0} A_{1} A_{2} A_{3} \right)^{-1} \\
    %
    \simeq & R_{0} {\rm Exp} \left( J_{l} \left( \lambda_{1} \boldsymbol \phi_{1} \right) \lambda_{1} \delta \boldsymbol \phi_{1} \right) {\rm Exp} \left( \lambda_{1} \boldsymbol \phi_{1} \right) A_{2} A_{3} \left( R_{0} A_{1} A_{2} A_{3} \right)^{-1} \\
    %
    = & R_{0} {\rm Exp} \left( J_{l} \left( \lambda_{1} \boldsymbol \phi_{1} \right) \lambda_{1} \delta \boldsymbol \phi_{1} \right) A_{1} A_{2} A_{3} A_{3}^{-1} A_{2}^{-1} A_{1}^{-1} R_{0}^{-1} \\
    %
    = & R_{0} {\rm Exp} \left( J_{l} \left( \lambda_{1} \boldsymbol \phi_{1} \right) \lambda_{1} \delta \boldsymbol \phi_{1} \right)  R_{0}^{-1} \\
    %
    = & {\rm Exp} \left( {\rm Ad}_{R_{0}} J_{l} \left( \lambda_{1} \boldsymbol \phi_{1} \right) \lambda_{1} \delta \boldsymbol \phi_{1} \right) \\
    = & I + \left( R_{0} J_{l} \left( \lambda_{1} \boldsymbol \phi_{1} \right) \lambda_{1} \delta \boldsymbol \phi_{1} \right)^{\wedge} + O \left( \left\| \delta \boldsymbol \phi_{1} \right\|^{2} \right)
  \end{split}
\end{align}
%
ここで$SO(3)$における性質である${\rm Ad}_{R} = R$を用いた．
同様に$j=2,3$の場合は以下の計算を行う．
%
\begin{align}
  \begin{split}
    & R_{0} A_{1} {\rm Exp} \left( \lambda_{2} \left( \boldsymbol \phi_{2} + \delta \boldsymbol \phi_{2} \right) \right) A_{3} \left( R_{0} A_{1} A_{2} A_{3} \right)^{-1} \\
    %
    \simeq & R_{0} A_{1} {\rm Exp} \left( J_{l} \left( \lambda_{2} \boldsymbol \phi_{2} \right) \lambda_{2} \delta \boldsymbol \phi_{2} \right) A_{1}^{-1} R_{0}^{-1} \\
    %
    = & {\rm Exp} \left( {\rm Ad}_{R_{0} A_{1} } J_{l} \left( \lambda_{2} \boldsymbol \phi_{2} \right) \lambda_{2} \delta \boldsymbol \phi_{2} \right) \\
    = & I + \left( R_{0} A_{1} J_{l} \left( \lambda_{2} \boldsymbol \phi_{2} \right) \lambda_{2} \delta \boldsymbol \phi_{2} \right)^{\wedge} + O \left( \left\| \delta \boldsymbol \phi_{2} \right\|^{2} \right)
  \end{split}
\end{align}
%
\begin{align}
  \begin{split}
    & R_{0} A_{1} A_{2} {\rm Exp} \left( \lambda_{3} \left( \boldsymbol \phi_{3} + \delta \boldsymbol \phi_{3} \right) \right) \left( R_{0} A_{1} A_{2} A_{3} \right)^{-1} \\
    %
    \simeq & R_{0} A_{1} A_{2} {\rm Exp} \left( J_{l} \left( \lambda_{3} \boldsymbol \phi_{3} \right) \lambda_{3} \delta \boldsymbol \phi_{3} \right) A_{2}^{-1} A_{1}^{-1} R_{0}^{-1} \\
    %
    = & {\rm Exp} \left( {\rm Ad}_{R_{0} A_{1} A_{2} } J_{l} \left( \lambda_{3} \boldsymbol \phi_{3} \right) \lambda_{3} \delta \boldsymbol \phi_{3} \right) \\
    = & I + \left( R_{0} A_{1} A_{2} J_{l} \left( \lambda_{3} \boldsymbol \phi_{3} \right) \lambda_{3} \delta \boldsymbol \phi_{3} \right)^{\wedge} + O \left( \left\| \delta \boldsymbol \phi_{3} \right\|^{2} \right)
  \end{split}
\end{align}
%
よって，$\partial R(t) / \partial \boldsymbol \phi_{j}$はそれぞれ以下となる．
%
\begin{align}
  \begin{gathered}
    \frac{ \partial R(t) }{ \partial \boldsymbol \phi_{1} }
    =
    R_{0} J_{l} \left( \lambda_{1} \boldsymbol \phi_{1} \right) \lambda_{1} \\
    %
    \frac{ \partial R(t) }{ \partial \boldsymbol \phi_{2} }
    =
    R_{0} A_{1} J_{l} \left( \lambda_{2} \boldsymbol \phi_{2} \right) \lambda_{2} \\
    %
    \frac{ \partial R(t) }{ \partial \boldsymbol \phi_{3} }
    =
    R_{0} A_{1} A_{2} J_{l} \left( \lambda_{3} \boldsymbol \phi_{3} \right) \lambda_{3} \\
  \end{gathered}
\end{align}
%

次に$\partial \boldsymbol \phi_{j} / \partial R_{j}$と$\partial \boldsymbol \phi_{j+1} / \partial R_{j}$に関してであるが，これは\ref{sec:smoothness}章で述べた回転に関するヤコビアンの計算と同じとなるため，詳細は省くが，計算結果はそれぞれ以下となる．
%
\begin{align}
  \frac{ \partial \boldsymbol \phi_{j} }{ \partial R_{j} }
  =
  \frac{ \partial {\rm Log} \left( \Delta R_{j} \right) }{ \partial \Delta R_{j} }
  \frac{ \partial \Delta R_{j} }{ \partial R_{j} }
  =
  J_{l}^{-1} \left( \Delta {\bf r}_{j} \right) R_{j-1}^{\top}
\end{align}
%
\begin{align}
  \frac{ \partial \boldsymbol \phi_{j+1} }{ \partial R_{j} }
  =
  \frac{ \partial {\rm Log} \left( \Delta R_{j+1} \right) }{ \partial \Delta R_{j+1} }
  \frac{ \partial \Delta R_{j+1} }{ \partial R_{j} }
  =
  -J_{l}^{-1} \left( \Delta {\bf r}_{j+1} \right) R_{j}^{\top}
\end{align}
%
ただし$\Delta R_{j} = R_{j-1}^{-1} R_{j}$，$\Delta {\bf r}_{j} = {\rm Log} \left( \Delta R_{j} \right)$とした．

これらから，$\partial R(t) / \partial R_{j}$はそれぞれ以下となる．
%
\begin{align}
  \begin{gathered}
    \frac{ \partial R(t) }{ \partial R_{0} }
    =
    I - R_{0} J_{l}^{-1} \left( \lambda_{1} \boldsymbol \phi_{1} \right) \lambda_{1} J_{l}^{-1} \left( \Delta {\bf r}_{1} \right) R_{0}^{\top} \\
    %
    \frac{ \partial R(t) }{ \partial R_{1} }
    =
    R_{0} J_{l}^{-1} \left( \lambda_{1} \boldsymbol \phi_{1} \right) \lambda_{1} J_{l}^{-1} \left( \Delta {\bf r}_{1} \right) R_{0}^{\top}
    -R_{0} A_{1} J_{l}^{-1} \left( \lambda_{2} \boldsymbol \phi_{2} \right) \lambda_{2} J_{l}^{-1} \left( \Delta {\bf r}_{2} \right) R_{1}^{\top} \\
    %
    \frac{ \partial R(t) }{ \partial R_{2} }
    =
    R_{0} A_{1} J_{l}^{-1} \left( \lambda_{2} \boldsymbol \phi_{2} \right) \lambda_{2} J_{l}^{-1} \left( \Delta {\bf r}_{2} \right) R_{1}^{\top}
    %
    -R_{0} A_{1} A_{2} J_{l}^{-1} \left( \lambda_{3} \boldsymbol \phi_{3} \right) \lambda_{3} J_{l}^{-1} \left( \Delta {\bf r}_{3} \right) R_{2}^{\top} \\
    %
    \frac{ \partial R(t) }{ \partial R_{3} }
    =
    R_{0} A_{1} A_{2} J_{l}^{-1} \left( \lambda_{3} \boldsymbol \phi_{3} \right) \lambda_{3} J_{l}^{-1} \left( \Delta {\bf r}_{3} \right) R_{2}^{\top}
  \end{gathered}
\end{align}
%
これらに$\partial {}^{L}r(t) / \partial R(t) = {\bf n}^{\top} \left( R(t) {}^{I}{\bf p} \right)^{\wedge}$を左から掛けた行列が$\partial {}^{L}r(t) / \partial R_{j}$となる．

